\chapter{Evaluation}
\label{chap:evaluation}

This chapter describes the experimental setup used to evaluate the optimisation approaches presented in \cref{chap:design-and-implementation} and presents the results of the performance evaluation. \Cref{sec:build-configurations} describes the build configurations used to isolate the effect of each optimisation flag. \Cref{sec:performance-evaluation} details the benchmark design, execution environment, and measurement methodology. Finally, \cref{sec:results-presentation} presents the collected performance data.

\section{Build Configurations}
\label{sec:build-configurations}

The three optimisation approaches are controlled by compile-time preprocessor flags: \texttt{ZUseOptimisedClearPath} for the optimised clear path, \texttt{ZUseSeparateDiscoveredLists} for the separate discovered list, and \texttt{ZUseDynamicArray} for the dynamic array. Each flag can be independently enabled (1) or disabled (0) via compiler flags passed during the JDK build.

Eight build configurations were created to evaluate all combinations of the three optimisation flags, as shown in \cref{tab:build-configs}. Each configuration produces a separate JDK build with its own binary, allowing direct comparison under identical benchmark conditions.

\begin{table}[htbp]
\centering
\caption{Build configurations and their enabled optimisation flags}
\label{tab:build-configs}
\begin{tabularx}{\textwidth}{lccX}
\toprule
\textbf{Configuration} & \textbf{Clear\ Path} & \textbf{Sep.\ Lists} & \textbf{Dyn.\ Array} \\
\midrule
\texttt{none}      & \textemdash & \textemdash & \textemdash \\
\texttt{dyn\_only} & \textemdash  & \textemdash & \checkmark \\
\texttt{sep\_only} & \textemdash & \checkmark  & \textemdash \\
\texttt{sep\_dyn}  & \textemdash & \checkmark  & \checkmark  \\
\texttt{clear\_path\_only} & \checkmark  & \textemdash & \textemdash \\
\texttt{clear\_path\_dyn}  & \checkmark & \textemdash  & \checkmark  \\
\texttt{clear\_path\_sep}  & \checkmark  & \checkmark  & \textemdash \\
\texttt{all}       & \checkmark  & \checkmark  & \checkmark  \\
\bottomrule
\end{tabularx}
\end{table}

Weak fields were evaluated as a separate configuration in the same benchmark harness. In this setup, the \texttt{WeakReference} benchmark class is replaced by its weak-field counterpart while using the baseline (\texttt{none}) ZGC optimisation flag configuration. This isolates representation-level effects from the three \texttt{WeakReference} pipeline optimisations, whilst still allowing direct comparison within the same execution and analysis workflow.

\section{Performance Evaluation}
\label{sec:performance-evaluation}

To evaluate the performance impact of the optimisation approaches, four microbenchmarks were used under controlled conditions: two \texttt{WeakReference}-based benchmarks and two weak-field counterparts with analogous workload structure. The benchmarks and their execution environment are described below.

\subsection{Benchmark Design}
\label{subsec:benchmark-design}

Two benchmark families were used to evaluate the optimisation approaches. The first family isolates reference-processing cost around a single shared referent, and the second is a more realistic multi-object workload that introduces greater allocation and liveness variability. Each family was implemented both with \texttt{WeakReference} objects and with weak fields. In order to enable \texttt{WeakReference}-to-weak-field comparison, both benchmark families use holder objects that either contain a \texttt{WeakReference} or an annotated weak field.

\subsubsection{Single-Object Benchmark}
\label{subsubsec:single-object-benchmark}

The Single-Object benchmarks allocate 10 million holder objects with either a \texttt{WeakReference} object or an annotated weak field each that all point to a single target object. The holder objects are allocated and stored in a randomised order using a Fisher-Yates shuffle~\cite{knuth97-taocp2} to avoid artificial locality patterns. None of the weak references are created with a \texttt{ReferenceQueue}. After a brief holding period, the single strong reference to the target object is released and a GC cycle is triggered manually. Because all weak references become clearable simultaneously in a single GC cycle, this benchmark maximises the per-cycle reference processing load and isolates the cost of the processing pipeline itself.

\subsubsection{Multi-Object Benchmark}
\label{subsubsec:multi-object-benchmark}

The multi-object benchmarks are a slightly more realistic workload. It allocates approximately 7.3 million objects, each containing a randomly sized byte array payload (between 509 and 1097 bytes). Each object has a strong reference and a corresponding holder object with either a \texttt{WeakReference} object or an annotated weak field. 5\% of the \texttt{WeakReference}s are created with a \texttt{ReferenceQueue} and 95\% without, reflecting a workload dominated by queue-less weak references but where the enqueue path is still exercised for a minority of references. The allocation and storing order of both strong references and holder objects are randomised using separate Fisher-Yates shuffles~\cite{knuth97-taocp2} to avoid artificial locality patterns.

Unlike the Single-Object benchmark, the strong references are not all released at once. Instead, the benchmark proceeds through five rounds of clearing. In each round, 20\% of the strong references are nulled in a randomised order, making the corresponding objects eligible for garbage collection. Between rounds, the benchmark sleeps to allow GC cycles to discover and process the weak references whose referent now is unreachable. This gradual release pattern means that weak references are discovered and processed across multiple GC cycles rather than in a single burst.

The benchmark runs three iterations with a five-second cooldown between them to allow the JVM to stabilise. Each iteration performs the full allocation, clearing, and GC cycle sequence independently.

\subsection{Execution Environment}
\label{subsec:execution-environment}

To ensure reproducible measurements, the benchmarks were executed under controlled conditions using a dedicated runner script. The script applies the following isolation measures:

\begin{itemize}
    \item \textbf{CPU pinning}: The Java process was pinned to a fixed set of CPU cores using \texttt{taskset}. Meanwhile, the monitoring and scripting processes ran on separate cores to avoid interference.
    \item \textbf{CPU frequency control}: The CPU governor is set to performance mode at the base clock frequency. Both the minimum and maximum scaling frequencies are locked to the base frequency to prevent frequency scaling from introducing variance.
    \item \textbf{Process priority}: The Java process is run with elevated priority (\texttt{nice -n -5}) to reduce scheduling interference from other system processes.
    \item \textbf{Cache clearing}: Filesystem caches are dropped between runs by writing to \texttt{/proc/sys/vm/drop\_caches}.
    \item \textbf{Cooldown period}: A configurable cooldown period (default 10 seconds) is inserted between variant runs to allow the system to return to a stable state.
\end{itemize}

For the single-object benchmark, the JVM was configured with a fixed 7\,GB heap and ZGC enabled. Tenuring thresholds are set to 1 to promote the target object to the old generation quickly, ensuring that the reference processor discovers it. GC statistics and reference processing logs were enabled to collect the results from.

For the multi-object benchmark the JVM was configured with a fixed 8\,GB heap, ZGC enabled, and a forced major collection interval of 0.5 seconds to ensure frequent GC activity during the benchmark. Tenuring thresholds are set to 1 to promote objects to the old generation quickly, ensuring that the reference processor discovers them. This is necessary since ZGC doesn't perform reference processing in the young generation. GC statistics and reference processing logs were enabled to collect the results from.

\subsection{Measurement Methodology}
\label{subsec:measurement-methodology}

Each benchmark was run for 100 outer iterations. Within each outer iteration, all eight \texttt{WeakReference} variants are run in sequence with environment preparation between each. In the same outer iteration, the weak-field benchmark variant is then executed as an additional run. This interleaved ordering mitigates systematic drift in system conditions affecting one configuration more than others.

During each run, a memory monitoring script samples memory metrics from the process at high frequency. Additionally, the current benchmark phase is tracked and for the multi-object benchmark the inner iteration is tracked as well. This tracking is done from the benchmark's log file to enable phase-aligned analysis. See \cref{tab:memory-monitoring-metrics} for the memory metrics collected.

In addition to the continuous memory samples, GC statistics are parsed from the log files of each run. For each GC metric, the parser extracts the Total Avg/Max pair and unit reported by ZGC. For the single-object benchmark, Max is used as the primary value for comparison; for the multi-object benchmark, Avg is used. The analysis pipeline also recognises the weak-field run naming convention and includes weak-field results as a dedicated \texttt{weak\_fields} configuration in comparative plots and summary tables. See \cref{tab:gc-log-metrics} for the GC log metrics collected.

\begin{table}[t]
\centering
\caption[Continuous monitoring metrics in the measurement pipeline]{Continuous monitoring metrics used in the measurement pipeline.}
\label{tab:memory-monitoring-metrics}
\begin{tabularx}{\textwidth}{@{}XlX@{}}
\toprule
\textbf{Metric} & \textbf{Units} & \textbf{Description} \\
\midrule
Timestamp & ms & Sample timestamp used to align and aggregate runs over time. \\
RSS & kB & Process resident set size (RSS), representing total physical memory in use by the JVM process. \\
GC Committed Memory & kB & GC-related memory currently committed by the JVM (actually backed memory from Native Memory Tracking). \\
Java Heap & kB & Java heap usage sampled from \texttt{jcmd GC.heap\_info}. \\
Phase & label & Benchmark phase label emitted by the benchmark and used for phase-aligned analysis. \\
Iteration & label & Inner-iteration identifier (multi-object benchmark) used to align repeated phase sequences. \\
\bottomrule
\end{tabularx}
\end{table}

\begin{table}[t]
\centering
\caption[GC log metrics in the measurement pipeline]{GC log metrics used in the measurement pipeline.}
\label{tab:gc-log-metrics}
\begin{tabularx}{\textwidth}{@{}XlX@{}}
\toprule
\textbf{Metric} & \textbf{Units} & \textbf{Description} \\
\midrule
Major Collection & ms & End-to-end duration of major collections reported by GC statistics. \\
Old Generation & ms & Total time spent in old-generation collection work. \\
Old Subphase: Concurrent Mark Follow & ms & Time spent in old-generation concurrent mark-follow processing. \\
Old Subphase: Concurrent References Process & ms & Time spent in concurrent reference processing in old-generation collection. \\
Young Subphase: Concurrent Mark Follow & ms & Time spent in young-generation concurrent mark-follow processing. \\
\bottomrule
\end{tabularx}
\end{table}

\section{Presentation of Results}
\label{sec:results-presentation}

This section presents the empirical results of the performance evaluation. The comparison covers the baseline (\texttt{none}) and all seven non-empty combinations of the three optimisation flags (clear path, separate lists, dynamic array) for \texttt{WeakReference}-based processing. In addition, a dedicated \texttt{weak\_fields} configuration is included, where the weak-field benchmark classes are executed in the same harness and analysed with the same tooling. This configuration set enables two complementary comparisons:
\begin{itemize}
	\item \textbf{Intra-model comparison} among \texttt{WeakReference} variants to isolate the impact of each optimisation and their interactions.
	\item \textbf{Inter-model comparison} between optimised \texttt{WeakReference} processing and weak-field representation to evaluate structural overhead trade-offs.
\end{itemize}

The primary performance metrics are the GC timing statistics extracted from ZGC logs, including major collection time, old-generation collection time, concurrent mark-follow subphases, and concurrent reference-processing subphases. Memory behaviour is characterised using continuous samples of RSS, GC committed memory, and Java heap usage. Following the methodology in \cref{subsec:measurement-methodology}, single-object benchmark comparisons use the per-run maximum values as primary statistics, while multi-object comparisons use per-run averages. This choice reflects the fact that single-object runs are designed to create concentrated peak load in one collection event, whereas multi-object runs distribute work over repeated phases.

\Cref{fig:metric-boxplots-single,fig:metric-boxplots-multi} show box plots of the GC timing metrics for the single-object and multi-object benchmarks. They summarise the distribution of per-run maxima for the single-object benchmark and per-run averages for the multi-object benchmark, allowing comparison of both central tendency and variability. \Cref{fig:continuous-memory-single,fig:continuous-memory-multi} show representative continuous-memory traces for the same benchmark families, illustrating how memory usage evolves over time and across GC phases.

\begin{figure}
	\centering
	\includegraphics[width=\textwidth]{metric_boxplots_single_1.pdf}
	\caption[GC metric box plots (single-object)]{Box plots of GC metrics across evaluated configurations for the single-object benchmark.}
	\label{fig:metric-boxplots-single}
\end{figure}

\begin{figure}
	\centering
	\includegraphics[width=\textwidth]{metric_boxplots_multi_2.pdf}
	\caption[GC metric box plots (multi-object)]{Box plots of GC metrics across evaluated configurations for the multi-object benchmark.}
	\label{fig:metric-boxplots-multi}
\end{figure}

\begin{figure}
	\centering
	\includegraphics[width=\textwidth]{continuous_memory_single_1.pdf}
	\caption[Continuous memory trace (single-object)]{Continuous memory trace for the single-object benchmark.}
	\label{fig:continuous-memory-single}
\end{figure}

\begin{figure}
	\centering
	\includegraphics[width=\textwidth]{continuous_memory_multi_2.pdf}
	\caption[Continuous memory trace (multi-object)]{Continuous memory trace for the multi-object benchmark.}
	\label{fig:continuous-memory-multi}
\end{figure}
